\section{Axiom 2: Conditional Independence of Block Error Rates}
\label{sec:axiom-2}

% Cross-references that will resolve in later plans:
%   \Cref{sec:model-order} -- chapter 2, drafted in Plan 2
%   \Cref{sec:bsc}         -- chapter 2, drafted in Plan 2
% Cross-references that will resolve in this plan (Task 5):
%   \Cref{sec:kronecker}   -- section 4.4, drafted in Task 5
% Cross-references that already resolve:
%   \Cref{sec:why-axioms}  -- section 4.1 (this chapter)
%   \Cref{sec:axiom-1}     -- section 4.2 (this chapter)
% Build warnings for the first two labels are expected during ch 4 drafting
% and will be resolved by Plan 4's Part I integration check.

\Cref{sec:axiom-1} gave the first independence axiom: error events at
distinct inputs are mutually independent, given the objective set and
the error rates. That axiom governs the \emph{spatial} structure of
errors, the relationship between different elements of the universal
set queried simultaneously or in any order.

This section gives the second axiom. Axiom 2 governs the \emph{stability}
of error rates across queries. It says that the error behavior on one
query tells you nothing about the error behavior on another query to the
same block: the per-block error rates are constant, not history-dependent.
The two axioms are logically independent of each other. Axiom 1 alone
does not imply Axiom 2, and Axiom 2 alone does not imply Axiom 1. Both
are needed, and their conjunction is exactly what makes the Kronecker
factorization theorem of \Cref{sec:kronecker} possible.

\subsection*{Partition blocks and model order}

Before stating Axiom 2, we need the partition structure it ranges over.
A full treatment appears in \Cref{sec:model-order}; we summarize here
the minimum needed for the axiom statement.

Given an objective set $A \subseteq U$, the Bernoulli model partitions the
universal set $U$ into blocks $B_1, \ldots, B_n$ such that every element
within a block has the same error rate. The \emph{model order} is the
number of blocks $n$. The simplest nontrivial case is order 2: the
partition $\{B_1, B_2\} = \{A,\, A^c\}$ separates the positives from
the negatives and assigns each group its own error rate.

For the Boolean universe $U = \{\top, \bot\}$ with objective set
$A = \{\top\}$, model order 2 yields two blocks:
\begin{align*}
  B_1 &= \{\top\} = A, \\
  B_2 &= \{\bot\} = A^c.
\end{align*}
Elements of $B_1$ have false-negative rate $\fnrate$ (they test positive
correctly with probability $\tprate$, incorrectly with probability
$\fnrate$). Elements of $B_2$ have false-positive rate $\fprate$ (they
test negative correctly with probability $\tnrate$, incorrectly with
probability $\fprate$). The channel matrix $Q$ introduced in
\Cref{sec:bsc} captures exactly these two rates:
\[
  Q \;=\;
  \begin{pmatrix}
    \tprate & \fnrate \\
    \fprate & \tnrate
  \end{pmatrix},
\]
where row $i$ describes the output distribution for an element in block
$B_i$.

In general, an $n$-th order model has $n$ blocks and a corresponding
channel matrix with $n$ distinct rows, one per block. Axiom 2 is a
statement about those rows.

\subsection*{The formal statement}

\begin{axiom}[Conditional independence of block error rates]
\label{ax:block-rates-indep}
Let $\{\alpha_1, \ldots, \alpha_n\}$ be the random error rates for the
respective partition blocks $B_1, \ldots, B_n$, conditional on the
objective set $A$. The rates $\alpha_1, \ldots, \alpha_n$ are mutually
independent:
\[
  f(\alpha_1, \ldots, \alpha_n \Given A)
  \;=\;
  \prod_{i=1}^n f(\alpha_i \Given A).
\]
\end{axiom}

In plain language: knowing the error rate on positives tells you nothing
about the error rate on negatives, and vice versa. The block rates are
independent of each other given the objective set. This applies to all
$n$ blocks simultaneously: no block's rate gives information about any
other block's rate.

A few clarifications help prevent common misreadings.

\medskip
\noindent\textbf{Rates as random variables.} The rates $\alpha_i$ are
random in the formal sense: they are the \emph{expected} error
frequencies, which may themselves vary across different realizations of
the data structure or algorithm. A Bloom filter, for instance, has a
false-positive rate that depends on the specific set of inserted
elements and the hash function choices. Axiom 2 says that the
resulting FPR and FNR are independent across blocks.

\medskip
\noindent\textbf{Stability, not constancy.} Axiom 2 does not say the
rates are the same for every element within a block (that is the
identically-distributed clause of Axiom 1). It says the rates are
jointly independent across blocks. The distinction matters: a
system could have rates that vary from element to element within a block
while still satisfying Axiom 2. The axioms play complementary roles;
neither subsumes the other.

\medskip
\noindent\textbf{The order-2 specialization.} In the second-order model,
Axiom 2 reduces to a single statement: the random false-positive rate
$\fprate$ and the random false-negative rate $\fnrate$ are
conditionally independent given $A$. The joint density of the
two rates factors:
\[
  f(\fprate, \fnrate \Given A)
  \;=\;
  f(\fprate \Given A) \cdot f(\fnrate \Given A).
\]
This is the instance that arises in virtually every example in
chapters 1--3.

\subsection*{Consequence: one channel matrix per block, stable across queries}

Axiom 2 has a structural consequence for how we model repeated queries.
Suppose we query an input $x \in B_i$ multiple times: first at step $t$,
then again at step $t'$, and so on. Axiom 2 says the error rate on query
at step $t$ does not influence the error rate on query at step $t'$.
Combined with Axiom 1 (which says errors at distinct elements are
independent), this means:

\begin{proposition}[Stable channel matrices]
\label{prop:stable-Q}
Under \Cref{ax:block-rates-indep}, the channel matrix row corresponding
to block $B_i$ is the same for every query involving an element of $B_i$,
regardless of the order or history of prior queries. Formally, for any
sequence of queries $(x_1, x_2, \ldots)$ with $x_t \in B_{i(t)}$, the
conditional distribution
\[
  \Prob{\tilde{A}(x_t) = s \Given A,\; x_t \in B_{i(t)}}
\]
is the same function of $i(t)$ and $s$ for all $t$, independent of what
happened at steps $1, \ldots, t-1$.
\end{proposition}

\begin{proof}
Fix block $B_i$ and a sequence of queries. The distribution of the
output at step $t$ for $x_t \in B_i$ is determined by the rate $\alpha_i$
assigned to block $B_i$. By Axiom 2, $\alpha_i$ is independent of
$\alpha_j$ for $j \neq i$. By Axiom 1, the error events at distinct
elements are independent. Together, the conditional distribution of the
output at step $t$ is a function only of $\alpha_i$ and the block
membership of $x_t$, not of any prior query outcomes. Since this holds
for all $t$, the channel matrix row for block $B_i$ is stable across the
entire query sequence.
\end{proof}

This proposition is the precise sense in which ``there is one channel
matrix $Q$ per input-block-pair.'' The matrix entry $Q_{ij}$ gives the
probability that an element of block $B_i$ produces output $j$. That
probability does not drift, does not warm up, and does not depend on
how many prior queries have been processed. The channel is genuinely
stationary within each block.

The practical consequence is large. Without this stability, we would need
a separate channel description for each time step, and the joint
distribution over a sequence of queries would require specifying an
exponentially growing family of time-varying matrices. With Axiom 2, a
single matrix $Q$ suffices for the entire query sequence. This is what
makes the Kronecker factorization of \Cref{sec:kronecker} possible: the
$m$-query joint confusion matrix is the Kronecker product of $m$ copies
of the same $Q$, precisely because each query sees the same stable
channel.

\subsection*{Worked illustration: the model-order-2 Boolean channel}

Let us make Axiom 2's content concrete with the simplest nontrivial case,
which will recur throughout the book: the Boolean universe with model
order 2.

Set $U = \{\top, \bot\}$ and $A = \{\top\}$, giving two blocks:
$B_1 = \{\top\}$ (positives) and $B_2 = \{\bot\}$ (negatives). The
channel matrix is
\[
  Q \;=\;
  \begin{pmatrix}
    \tprate & \fnrate \\
    \fprate & \tnrate
  \end{pmatrix}
  \;=\;
  \begin{pmatrix}
    1 - \fnrate & \fnrate \\
    \fprate & 1 - \fprate
  \end{pmatrix}.
\]

Now consider a sequence of three queries: query 1 asks for $\top$,
query 2 asks for $\bot$, query 3 asks for $\top$ again. Axiom 2 says
the following.

\begin{itemize}
  \item Query 1 involves $\top \in B_1$. It produces a positive response
    with probability $\tprate$ and a negative response with probability
    $\fnrate$. These probabilities come from row 1 of $Q$.

  \item Query 2 involves $\bot \in B_2$. It produces a positive response
    with probability $\fprate$ and a negative response with probability
    $\tnrate$. These probabilities come from row 2 of $Q$. \emph{The
    outcome of query 1 is irrelevant.}

  \item Query 3 involves $\top \in B_1$ again. It again uses row 1 of
    $Q$, with the \emph{same} probabilities $\tprate$ and $\fnrate$ as
    query 1. The fact that query 1 may have made an error (false negative)
    does not change the probability of error on query 3.
\end{itemize}

This is the channel's stationarity within blocks: the same row of $Q$
applies whenever the same block is queried, no matter what came before.

To write the joint distribution over the three-query sequence
$(x_1, x_2, x_3) = (\top, \bot, \top)$ with outcomes
$(s_1, s_2, s_3)$, Axioms 1 and 2 together give:
\begin{align*}
  &\Prob{\tilde{A}(\top) = s_1,\; \tilde{A}(\bot) = s_2,\; \tilde{A}(\top) = s_3 \Given A} \\
  &\quad=\; Q_{1,s_1} \cdot Q_{2,s_2} \cdot Q_{1,s_3},
\end{align*}
where $Q_{i,s}$ denotes the $(i,s)$-entry of $Q$. The distribution
factors as a product of independent channel evaluations, each determined
by the block of the queried element. No interaction terms appear; no
history from query 1 flows into query 3.

For a concrete numerical example, take $\fprate = 0.05$ and
$\tprate = 0.90$ (so $\fnrate = 0.10$ and $\tnrate = 0.95$). The
probability that the three-query sequence produces outcomes
$(+, -, +)$, that is, two true positives and one true negative, is:
\[
  Q_{1,+} \cdot Q_{2,-} \cdot Q_{1,+}
  \;=\;
  0.90 \times 0.95 \times 0.90
  \;=\;
  0.7695.
\]
The probability that query 2 produces a false positive while queries 1
and 3 produce true positives is:
\[
  Q_{1,+} \cdot Q_{2,+} \cdot Q_{1,+}
  \;=\;
  0.90 \times 0.05 \times 0.90
  \;=\;
  0.0405.
\]
Both numbers follow from reading the appropriate rows of $Q$ and
multiplying. No additional modeling is needed. The stationarity
guaranteed by Axiom 2 is what makes this reading-and-multiplying
procedure valid.

\subsection*{Contrast: a Markov-correlated channel violates Axiom 2}

To see what Axiom 2 rules out, consider a channel that violates it
explicitly. A \emph{Markov-correlated channel} lets the error rate on
each query depend on the outcome of the previous query. Concretely:

\begin{enumerate}
  \item If the previous query produced a correct response, the current
    query uses a ``base'' false-positive rate $\fprate_0$.
  \item If the previous query produced an error, the current query uses
    an elevated false-positive rate $\fprate_1 > \fprate_0$ (errors
    cluster together).
\end{enumerate}

This is a plausible model for hardware memory with correlated faults, or
for a sensor that drifts after one bad reading. But it violates Axiom 2.

To see why, consider block $B_2 = \{\bot\}$ (negatives). On query $t$,
the error rate for block $B_2$ is not a fixed value $\fprate$: it is
$\fprate_0$ if the previous query was correct, and $\fprate_1$ if the
previous query was an error. The rate $\alpha_2$ at step $t$ is not
independent of the rate at step $t-1$; they are coupled through the
Markov chain. The joint density of the rates at different steps does not
factor, and Axiom 2's independence condition fails.

Let us put numbers to this. Suppose the base rate is $\fprate_0 = 0.05$
and the elevated rate is $\fprate_1 = 0.30$. The channel starts in the
correct state. Over two successive queries to $\bot$, the joint
distribution of outcomes is:

\begin{center}
\begin{tabular}{lccl}
\hline
Outcome pair & Probability & Axiom-2 product & Discrepancy \\
\hline
$(-, -)$ & $0.95 \times 0.95 = 0.9025$ & $0.9025$ & none \\
$(-, +)$ & $0.95 \times 0.05 = 0.0475$ & $0.0475$ & none \\
$(+, +)$ & $0.05 \times 0.30 = 0.0150$ & $0.0025$ & $6\times$ inflation \\
$(+, -)$ & $0.05 \times 0.70 = 0.0350$ & $0.0475$ & deflation \\
\hline
\end{tabular}
\end{center}

Here the marginal false-positive rate is $0.95 \times 0.05 + 0.05 \times 0.30 = 0.0625$
(averaging over both states), so the Axiom-2 product for two positives
would be $0.0625^2 = 0.0039$. The actual probability of two consecutive
false positives is $0.05 \times 0.30 = 0.0150$, a factor of $3.8\times$ higher.

The discrepancy is not small. The Markov channel produces error
clustering that is invisible to the factored model. If a practitioner
were to apply the Kronecker factorization of \Cref{sec:kronecker} to
this channel, they would underestimate the probability of back-to-back
errors by a factor of nearly four. In applications where error
clustering is costly, this is not a tolerable approximation error.

The Markov channel is not a pathological construction: it models a
real failure mode. What Axiom 2 asserts is that the Bernoulli framework,
in its standard form, is designed for channels where this clustering does
not occur. If the channel is known to be Markov-correlated, the model
must be extended, for example by conditioning on the hidden state, or
by treating the full trajectory as the object to be analyzed. The
Bernoulli axioms are then either violated, requiring a richer model, or
satisfied at the level of the extended state space, in which case the
standard machinery applies after the appropriate state augmentation.

\subsection*{What Axiom 2 earns us}

Axiom 1 and Axiom 2 are the full axiomatic foundation of the Bernoulli
framework. They are minimal in the sense that neither alone is sufficient
for the Kronecker factorization theorem: Axiom 1 gives spatial
independence across elements, Axiom 2 gives temporal stability of rates
across queries. Together they imply that the joint confusion matrix
over any sequence of queries is a Kronecker product of single-query
matrices, one factor per query (\Cref{sec:kronecker}).

The parametric savings are substantial. A general $n$-th order model
describing $m$ joint queries, without any independence assumption, has
$n(2^m - 1)$ free parameters: $2^m$ output symbols, $n$ rows in the
channel matrix, one free parameter per output symbol per row. Under
Axioms 1 and 2, the Kronecker structure forces the parameter count down
to $2m$ regardless of $n$: two parameters per query (one per block in
the order-2 model), with no cross-query parameters at all. For $m = 50$
queries on an order-2 channel, this is the difference between $50$
parameters (under the axioms) and $50 \times (2^{50} - 1) \approx 5.6
\times 10^{16}$ parameters (in the general case). The axioms are not
merely philosophically satisfying; they are what makes the model
tractable.

We prove the Kronecker factorization theorem and its corollary on
parameter reduction in \Cref{sec:kronecker}. That section is where the
full benefit of both axioms becomes visible.
